\section{Agent Prompt Templates}
\label{app:prompts}

This section presents the core prompt templates used in our generation pipeline. Each prompt is designed using a "Role-Task-Constraint" structure, injecting the structured context data defined in Appendix~\ref{app:data_structures}.

\subsection{Structural Planning Stage Prompts}
\label{app:plan_prompt}
The \textbf{Infer Global Skeleton Prompt} is responsible for generating the general skeleton. 
\label{app:plan_prompt_skeleton}
\begin{promptbox}[Infer Global Skeleton Prompt Template]
## [Role]
You are a **DEVS System Architect**. Your task is to design the overall module hierarchy for a DEVS simulation system.

## [Input]
**System Name**: `{root_name}`
**Requirements**:
{requirements}

## [Task]
Decompose the system into a hierarchical module structure. Return a list of modules.

## [Rules]
1. The first module MUST be the root system: `{root_name}`. It should be a coupled model with children.
2. Every name mentioned in any `children_names` MUST appear as a `name` somewhere in the list.
3. Use hierarchical decomposition: group related functionality into intermediate coupled models before reaching atomic models.
4. Keep descriptions SHORT (1-2 sentences). Only state what the module does.
5. Module names should be valid Python class names (PascalCase, no spaces/special chars).
6. Atomic (leaf) modules should have `children_names: []`.
7. Do NOT over-decompose. Less than 4 levels of hierarchy is usually sufficient.
8. Each coupled model should have at least 2 children (you can claim one type being instantiated multiple times)
9. DEVS Principles:
  - A system is hierarchical.
  - Atomic models have behavior (state machines) but NO sub-components.
  - Coupled models have sub-components and routing (couplings) but NO behavior.
  - The input/output ports of a model can only connect to a sibling (IC) or be a proxy of parent model (EIC, EOC).
  - Define the data_flows to establish the exact communication topology between the modules you create.

## [IMPORTANT: KEEP IT SIMPLE]
- **Minimize the number of modules**: Only create modules that are truly necessary.
- **Prefer shallow hierarchies**: A flat structure with fewer levels is better than a deep one.
- **Each module should have a clear, distinct responsibility**: Avoid overlapping functions.
- **Atomic modules should be self-contained**: Each should implement a complete, coherent piece of functionality.

## [Field Guidance]
- `name`: Valid Python identifier, PascalCase. Example: "InputHandler", "CoreProcessor".
- `description`: 1-2 sentences stating what the module does. Example: "Validates and normalizes incoming data.". Also describe the data flow inside the model (EIC, EOC, IC). 
- `children_names`: List of child module names. Empty list `[]` for atomic (leaf) modules.

## [Example]
For a system "DataPipeline" with requirements "ingest, transform, and output data":
- modules:
  - name: "DataPipeline", description: "Top-level coordinator. Routes data through ingestion, transformation, and output stages.", children_names: ["Ingester", "Transformer", "Emitter"]
  - name: "Ingester", description: "Reads raw data from source and validates format.", children_names: []
  - name: "Transformer", description: "Applies transformation rules to validate data.", children_names: ["Mapper", "Filter"]
  - name: "Mapper", description: "Maps input fields to target schema.", children_names: []
  - name: "Filter", description: "Removes invalid or duplicate records.", children_names: []
  - name: "Emitter", description: "Writes transformed data to target destination.", children_names: []
\end{promptbox}

Following the global skeleton inference, the architecture refinement process uses two distinct prompt templates depending on the model type being processed. This corresponds to the top-down traversal defined in Stage 1. 

The \textbf{Refine Coupled Plan Prompt} serves a dual purpose: it defines the internal routing (EIC, IC, EOC) for the current coupled model, and simultaneously generates the preliminary plan for all of its direct children.
\label{app:plan_prompt_coupled}
\begin{promptbox}[Refine Coupled Plan Prompt Template]
## [Role]
You are a **DEVS System Architect**. 

## [Input]
**Target Module**: `{target_name}`
**Children in Global Plan**: `{children_names_str}`
**Original Requirements**:
{requirements}

**Global Plan Overview** (full module hierarchy):
{global_plan_str}

**Model's Simple Plan** (this model's initial interface from parent):
{parent_simple_str}
**Parent's Detailed Plan** (system context):
{parent_detail_str}

## [Task]
Generate a detailed specification for the COUPLED module `{target_name}` and simple specifications for its direct children ({children_names_str}).

### STEP 1: Design Coupled Wrapper (`detailed_plan`)
- **function**: Describe the overall purpose and capability of this entire subsystem as a unified whole. While the coupled wrapper itself only contains structural connections (no active routing/state logic), this field should summarize what the encapsulated subsystem achieves.
- **model_init_args**: inherit from Parent's Simple Plan.
- **input_ports / output_ports**: inherit from Parent's Simple Plan.

### STEP 2: Design Children Contracts (`children_plans`)
- **function**: 1-2 sentences describing child's responsibility.
- **logging**: 1-2 sentences, mention what logging items should be covered.
- **model_init_args**: full args (name/type/structure); always start with name (str) and parent (Coupled | None).
- **input_ports / output_ports**: full port definitions for sibling and parent matching. MUST make sure the structures of ports match strictly.

### STEP 3: Design `coupling_specification`:
- MUST define the network routing here using ONLY these 3 strict DEVS patterns. List them clearly line-by-line:
    1. **EIC (External Input Coupling)**: `parent.IN.port_name -> child.IN.port_name`
    2. **IC (Internal Coupling)**: `child_A.OUT.port_name -> child_B.IN.port_name`
    3. **EOC (External Output Coupling)**: `child.OUT.port_name -> parent.OUT.port_name`
- **CRITICAL COUPLING RULES**:
    1. **NO HALLUCINATIONS**: Every port name you use MUST EXACTLY MATCH either the Parent's inherited ports or the Children's ports you defined in STEP 2.
    2. Not all ports need parent connections; children can communicate entirely via IC.

## [Field Guidance]
- For ANY `dict` or `list` types, you MUST use strict Python representation (e.g., `{'sequence_number': int, 'is_retry': bool}`).
- Port protocol: Must include initial_state (state at T=0), initial_signal (signal at startup), description.
\end{promptbox}

\vspace{1em}

The \textbf{Refine Atomic Plan Prompt} takes the preliminary simple plan assigned by its parent coupled model and expands it into a detailed plan, focusing strictly on state transitions, timing, and local logic without altering the assigned interfaces.
\label{app:plan_prompt_atomic}
\begin{promptbox}[Refine Atomic Plan Prompt Template]
## [Role]
You are a **DEVS System Architect**. 

## [Input]
**Target Module**: `{target_name}`
**Children in Global Plan**: None (leaf)
**Original Requirements**:
{requirements}

**Global Plan Overview** (full module hierarchy):
{global_plan_str}

**Model's Simple Plan** (this model's exact contract assigned by parent):
{parent_simple_str}
**Parent's Detailed Plan** (system context):
{parent_detail_str}

## [Task]
Generate a detailed specification for the ATOMIC module `{target_name}`.

### For the detailed ATOMIC `{target_name}` plan:
- **function**: Describe pure responsibility, state machine, event handling, and timing. Do NOT describe logging here.
- **logging**: Extract ALL logging/output requirements from the original requirements that apply to this model. Include payload structure, format, and timing.
- **model_init_args**: Strictly copy from Model's Simple Plan.
- **input_ports / output_ports**: Strictly copy from Model's Simple Plan. DO NOT invent new ports.

## [Field Guidance]
- For ANY `dict` or `list` types, you MUST use strict Python representation (e.g., `[{'job_id': int, 'priority': float}]`).
- Port protocol: Must include initial_state (state at T=0), initial_signal (signal at startup), description.
\end{promptbox}

\subsection{Behavioral Synthesizing Stage Prompts}
\label{app:code_prompt}
The \textbf{Model Creator} synthesizes the actual executable Python code. To ensure high-quality generation, we inject a set of global engineering standards alongside model-specific instructions (Atomic vs. Coupled).

The main prompt template is: 
\begin{promptbox}[Main Code Generation Prompt Template]
## [Task]
Construct a complete Python file containing a **{model_type} DEVS model** named `{name}` using `xdevs.py`.

{global_standards}

{model_specific_instructions}

{feedback}

## [Context Info]
**Sub-Models (for Coupled definitions)**: 
{sub_models}

**System Context**:
(The environment around this model)
(You must especially guarantee the JSONL output requirements are met.)
{context_str}

## [Utils]
{util_desc}

## [Class Definitions]
{definitions}

## [Specification]
The ports, logic, logging dict keys, and parameters of the model should strictly follow the specification (including their types, functions), only two can be added / modified: in __init__ args, `name: str`, and `parent: Coupled | None`:
{spec}

## [Reference Example]
Refer to this example for coding style and imports:
{example}

## [Output]
Return the Python code enclosed in <python_code> tags. 
Do not use markdown backticks.

Example:
Think step by step, decompose the requirements and state machine.
Finally the enclosed code.
<python_code>
class MyModel(Atomic):
    ...
</python_code>
\end{promptbox}

\begin{promptbox}[Global Engineering Standards (Injected into Main Prompt)]
## [Global Standards]
Adhere to the following engineering standards for all model types:

### 1. Imports \& Dependencies
- **Whitelist**: Restrict imports to the following packages: `numpy`, `math`, `random`, `time`, `pandas`, `xdevs` (and `xdevs.models`).
- **Project Utils**: Import necessary utilities (e.g., `get_sim_logger`, `get_current_time`) from `devs_project.devs_utils.xxx`. Refer to [Utils] for detailed import statements.
- Other submodels in the project can be imported as needed.

### 2. Data Protocol \& Typing
- You are only allowed to use the following Atomic Primitives and Composite Types in ports and arguments. 
- **Atomic Primitives**: Only `int`, `float`, `str`, `bool` instances for all base values.
- **Composite Types**: Only `dict` and `list`.
- **Consistency**: Ensure that all ports and arguments are the same as those stated in the [Specification]. 
- **Recursive Schema Definition**:
    - Recursively define the structure of all composite arguments (in `__init__`) and port data types (in Docstrings).
    - Continue the definition until all fields resolve to Atomic Primitives.
    - **Dictionaries**: Explicitly list every Key name, the Type of its Value, and the explanation of it. 
    - **Lists**: Explicitly state the Type of elements contained in the list.
- Special: only the `parent` argument of the `__init__` method can be of type `Coupled | None`.

### 3. Type Docstring Schema
- **Structure**: Follow this hierarchy for describing types in Class and Method docstrings:
    - Root: `name (type): Description`
    - Dict Keys: Indent 2-4 spaces, `key_name (type): Description`
    - List Items: Indent 2-4 spaces, `- (type): Description`
- *Exception*: For Logging, you can refer to the structure described in the ports section, as they are often the same.
- **Example**:
    ```python
    \"\"\"
    ...
    - in_packet (dict): Network packet.
        header (dict): Protocol header.
            src_ip (str): Source IP.
        payload (str): Data content.
    - batch_updates (list[dict]): Updates.
        - (dict): Single update.
            node_id (int): Target ID.
    \"\"\"
    ```

### 4. Coding Conventions
- **Explicit Configuration**: Define all configuration parameters explicitly in `__init__`. Omit `*args` and `**kwargs`.
- **Logging**: Log all key events using `self.logger.info({{"key": "value"}}, log_type=...)`. The main msg must be a dict. They must have the exact precise structure and keys as the context (especially the system goal) and the Model Specification. You should explain the structure of the data in the docstring. However, you only need to list the logging happen in this model, those in submodels are strictly forbiddened.
- **Parameter Storage**: Store internal hardcoded parameters (not passed via `__init__`) in a `self.param` dictionary.
\end{promptbox}

\begin{promptbox}[Atomic Model Instructions (Injected into Main Prompt)]
### [Atomic Model Specifics]
0. Must import: `from xdevs.models import Atomic, Coupled, Port`. Atomic for inherit, Coupled for __init__ arg type.
1. **Inheritance**: Inherit from `Atomic`.
2. **Docstring**: The class MUST include a standard class docstring strictly following this format:
    ```python
    class {name}(Atomic):
        \"\"\"
        Function: 
            - ...General function
            - ...Every state: how it transfer, and what to output after the state is over.
        Logging in this model:
            - ...
            - ...
        Input Ports:
          - port_name (type): description
            structure: ...
            protocol: initialize: ... ; process: ...
        Output Ports:
          - port_name (type): description
            structure: ...
            protocol: initialize: ... ; process: ...
        \"\"\"
    ```
3. **Constructor (`__init__`)**:
    - Signature: `def __init__(self, name: str, parent: Coupled | None, <explicit_config_args>)`
    - Docstring: should have a docstring describing the arguments, including the detailed type and description. using the following format:
        ```python
        \"\"\"
        Args:
            name (str): The unique name of the model.
            parent (Coupled | None): the parent model. If None, the model is a root model.
            arg_name1 (type): description
        \"\"\"
        ```
    - Steps:
        1. Call `super().__init__(name)`.
        2. Assign `self.parent = parent`.
        3. Initialize logger: `self.logger = get_sim_logger(self)`.
        4. Register Ports: Use `self.add_in_port(Port(type, "name"))` and `self.add_out_port(Port(type, "name"))`.
        5. Initialize State: Set member variables and call `self.hold_in(phase, time)`. 
        6. Log creation: `self.logger.info({{keys: values, ...}}, log_type=...)`
4. **Core Behaviors**:
    - Implement `initialize(self)`: Set initial state. Set phase/sigma using `self.hold_in(phase, time)`. Log initialization.
        - It can not send any output. If you need to send a initial signal (e.g. report you are ready), you can use `self.hold_in(phase, time)` to schedule the event, prepare the payload, and send it in `lambdaf`.
        - If any port has `initial_signal` to emit immediately or after a duration, the `initialize` method **MUST** schedule an output event using `self.hold_in("SOME_STATE", duration)`.
    - Implement `lambdaf(self)`: Only do the output, any other operations should be done in the following `deltint`:
        - Send output via `self.output["port"].add(payload)`.
        - DO NOT change the state, sigma, kpi_counter, etc. Leave that to the following `deltint`.
    - Implement `deltint(self)`: Only do the following:
        - Get the old internal state: `self.phase`. Get total time(from last state change to expected next state change, which is just the sigma set last time): `self.ta()`.
        - Handle internal timeouts. And update the internal queue / kpi_counter / etc. accordingly. (Because deltint is called right after lambdaf, which means the output of the old state is already sent). 
        - Prepare the payload of the output of the next phase. Make sure the prepared payload is the one used in `lambdaf` of the new phase.
        - Always schedule next internal event in the end: `self.hold_in(phase, sigma)`. If not interrupted, the model will emit the output prepared after sigma time units. 
        - Log events (if needed). 
    - Implement `deltext(self, e)`: Only do the following:
        - Handle external events (`self.input["port"].values`).
        - Get internal state: `self.phase`. Get total time(from last state change to expected next state change, which is just the sigma set last time): `self.ta()`.
        - Prepare the payload of the next lambdaf. Make sure the prepared payload variable is the one used in `lambdaf`.
        - Always schedule next internal event in the end: `self.hold_in(phase, sigma)`.
        - Log events (if needed). 
    - Implement `exit(self)`: Cleanup and final stats logging.
    - **Event Handling Logic**:
        - **Execution Sequence (CRITICAL)**: `lambdaf` will send outputs before `deltint` schedules the next internal event. Thus, the payload sent in `lambdaf` should be prepared in the previous `deltint`, `deltext`, or `initialize`. 
        - **Confluent Events (`deltcon`)**: By default, internal events (`deltint`) take precedence over external events when they occur simultaneously. Explicitly override the `deltcon(self)` method ONLY IF you need to change this logic (e.g., to process external events first).
        - **Initialization**: Realize the ports.protocol's initialize descriptions: 
            - initial_signal: If a signal or information should be sent at initialization(i.e. protocol.initial_signal), you can use `self.hold_in("INIT", 0)` to schedule the event and send it in `lambdaf`. This is the only way to send a signal at initialization.
            - initial_state: modify the logic and initial values to make sure it is realized. 
5. **logging requirements**: make sure all the events required are logged. And the keys and structures of the logs must match the Specification exactly. 
\end{promptbox}

\begin{promptbox}[Coupled Model Instructions (Injected into Main Prompt)]
### [Coupled Model Specifics]
0. Must import: `from xdevs.models import Atomic, Coupled, Port`
1. **Inheritance**: Inherit from `xdevs.models.Coupled`.
2. **Docstring**: The class MUST include a standard class docstring strictly following this format:
    ```python
    class {name}(Coupled):
        \"\"\"
        Function: 
          - ...
          - ...
          - Sub-models: 
            - sub_model_class_name: name=sub_model_instance_name. Brief description.
        Logging in this model:
          - ...
          - ...
        Input Ports:
          - port_name (type): description
            structure: ...
            protocol: initialize: ... ; process: ...
        Output Ports:
          - port_name (type): description
            structure: ...
            protocol: initialize: ... ; process: ...
        \"\"\"
    ```
3. **Container Logic**: Treat this class as a pure structure container. Implement ONLY `__init__`.
4. **Sub-models Imports**: Use relative imports for sub-models (e.g., `from .folder.file import SubModelName`).
5. **Constructor (`__init__`)**:
    - Signature: `def __init__(self, name: str, parent: Coupled | None, <explicit_config_args>)`
    - Docstring: should have a docstring describing the arguments, including the detailed type and description. using the following format:
        ```python
        \"\"\"
        Args:
            name (str): The unique name of the model.
            parent (Coupled | None): the parent model. If None, the model is a root model.
            arg_name1 (type): description
        \"\"\"
        ```
    - Steps:
        1. Call `super().__init__(name)`.
        2. Assign `self.parent = parent`.
        3. Initialize logger: `self.logger = get_sim_logger(self)`.
        4. Register Ports: Use `self.add_in_port(...)` and `self.add_out_port(...)`.
        5. Instantiate Components: Create sub-model instances and register them via `self.add_component(instance)`.
        6. Define Couplings: Use `self.add_coupling(src, dst)` for:
            - **EIC**: `self.input["port_name"]` -> `sub.input["port_name"]`
            - **IC**: `sub_a.output["port_name"]` -> `sub_b.input["port_name"]`
            - **EOC**: `sub.output["port_name"]` -> `self.output["port_name"]`
        7. Log creation: `self.logger.info(...)`
    - Note: For steps 5-6, you should refer to Sub-Models to get the right init args names and port names. These information can be used as a correction and supplement to the coupling logic (in case some names are inconsistent). 
\end{promptbox}


The \textbf{Model Summarizer} implements the bottom-up adaptation mechanism. It analyzes the generated code to extract the "Ground Truth" interface, which is then used to guide the generation of the parent coupled model.

\begin{promptbox}[Summarizer Prompt Template]
## [Task]
Analyze the provided Python code for a DEVS model.
Extract the model's metadata into a strict structure based on the schema below. 

## [Rules]
- **model_init_args**: You must extract all the arguements from the `__init__` method (except `self`). You should make sure that the usage of all args are clear (especially for the str, dict types, and possible options).
- **logging**: Look for the docstring and logging usage to fill `logging`. You must clearly state the `log_type` and the structure of main msg. 
  - You only need to extract the logger used in this model. The logging inside sub-models are not needed, as they are handled by their own models.
  - If it is described in the docstring, just copy the releted part in docstring, you can add more details if needed. (e.g. to keep the name of data clear, or other needs. )
- **ports**: Refer to the docstring for port definitions. 
- **function**: Look for the main logic of the model. For communication protocols, find how it starts, send, receive, and ends(if any). 
  - For coupled models, you must copy the submodels described in the docstring, especially the relation between the class name and the instance name.

{feedback}

## [Sub-Models Info]
{sub_models}

## [Code]
```python
{code}
```
\end{promptbox}

\subsection{Godot Generation Agent Prompt}

We use the following prompt for the Godot construction. 

\begin{promptbox}[Godot Generation Agent Prompt Template]
You are a coding agent using godogen to build one runnable Godot scene for a DEVS simulator.

Target engine:

- Godot 4.6 (STRICT). All generated code and APIs must be compatible with Godot 4.6 only.

Available inputs:
1) DEVS_examples (reference behaviors/patterns),
2) one scenario YAML file,
3) one DEVS simulator package (if provided in workspace),
4) image generation API config below.

Primary goal:
Produce a Godot project/scene that can be opened and run directly in Godot 4.6, and that visualizes the DEVS scenario behavior correctly.

Hard requirements:
1) Build scene logic from YAML + DEVS_examples (+ simulator package if available).
2) Connect runtime events via MQTT (or provided runtime stream) and update scene deterministically.
3) Generate and APPLY textures for scene elements (do not leave plain placeholders):

   - background image
   - source/facility visual
   - destination visual
   - moving entities (e.g., aircraft/vehicles/agents)
   - cargo/package/queue-related visuals
   - core UI icons for status panels
     4) Ensure all major visual elements are textured in final scene.
     5) Prioritize delivering runnable output over long explanations.

Layout and visibility requirements:
6) Scene must be fully visible and readable at:

   - 1152x648
   - 1366x768
   - 1920x1080
     7) Use responsive UI anchors and scalable center viewport.
     8) Avoid clipping critical panels/text and avoid overlap between key UI and moving entities.

Runtime motion requirements:
9) Scene must show visible motion after launch.
10) If live MQTT/runtime stream is available, drive updates from incoming events.
11) If live stream is unavailable within 3 seconds, switch to replay/demo mode (same event schema) so the scene is still animated and demonstrable.

Godot 4.6 compatibility rules (STRICT):
12) Do not use Godot 3-only APIs.
13) For UDP in Godot 4.6:

   - use PacketPeerUDP.new()
   - bind to local port
   - poll packets in _process(delta) using get_available_packet_count()
   - read via get_packet()
   - parse payload and apply updates
     14) NEVER call nonexistent methods such as:
   - PacketPeerUDP.set_blocking_mode(...)
     15) On bind/read/parse failure: log error and continue safely without crashing.

GDScript type-safety rules (STRICT):
16) Assume warnings may be treated as errors. Generated code must be warning-free.
17) Do NOT rely on inferred typing when RHS may be Variant.
18) For values from JSON/Dictionary/Array/get(), always use explicit type annotations plus cast/conversion.
19) Avoid := when RHS can be Variant; prefer explicit declarations.
20) Always convert Dictionary.get(...) values before assignment/use.
21) Before final output, fix all type warnings.

Type-safe style examples:

- var msg: Dictionary = JSON.parse_string(text) as Dictionary
- var payload: Dictionary = msg.get("payload", {}) as Dictionary
- var event_type: String = str(msg.get("event_type", ""))
- var ts: float = float(msg.get("timestamp", 0.0))
- var records: Array = msg.get("records", []) as Array

Execution policy:

- Do not spend time on extended self-check reports.
- Do not iterate endlessly on optimization.
- Finish when scene is runnable, textured, visibly animated, and warning-free.
- If one asset generation fails, retry briefly, then generate a substitute asset and continue.

Final output:
A) Runnable Godot scene/project files (Godot 4.6 compatible)
B) Brief run note: which scene to open and run
C) Brief list of generated texture files actually used in scene
D) Confirmation that no warnings remain (especially Variant-inference warnings)
\end{promptbox}